% GENERATED FILE. Edit canonical lessons, then run npm run generate:books.
% canonical-source-hash: fnv1a64:b9a7a141919c12eb
% canonical-lessons: JA-C193-jinja, JA-C193-tera, JA-C193-eigakan, JA-C193-bijutsukan, JA-C193-mieru

\chapter{Shrines, Temples and Cinemas}
\label{ch:ja-shrines-temples-and-cinemas}
\hlchaptermodality{193}

\begin{chapteropening}
\textbf{By the end of this chapter:} \emph{I can name the sights of a town.}

Complete the last lesson of chapter 193: I can name the sights of a town.
\end{chapteropening}

\section[\texorpdfstring{\ja{じんじゃ} (jinja)}{じんじゃ (jinja)}]{\ja{じんじゃ} (jinja) — shrine; Shinto shrine}

\label{lesson:JA-C193-jinja}



\begin{quote}
Before the new one: say the Japanese for to ask, then the Japanese for guidance.
\end{quote}

\subsection*{You'll want to know: \ja{じんじゃ}}
\textbf{\ja{じんじゃ}} — \emph{jinja} — \textquotedblleft{}shrine; Shinto shrine\textquotedblright{}.

\begin{quote}
\textbf{\ja{まいとし} \ja{じんじゃに} \ja{いきます。}} — \emph{maitoshi jinja ni ikimasu} — \textquotedblleft{}I go to the shrine every year.\textquotedblright{}
\end{quote}

\begin{cousinweb}
\textbf{\ja{じんじゃ}} (\emph{jinja}) is Sino-Japanese (kango): two Chinese characters read with their on'yomi, readings borrowed from Chinese. They are \emph{jin}, \textquotedblleft{}god, spirit\textquotedblright{} (Middle Chinese \emph{zyin}) and \emph{ja}, \textquotedblleft{}shrine, society\textquotedblright{} (\emph{dzyaeX}) (KANJIDIC2; Middle Chinese in Baxter's transcription).
\end{cousinweb}

\begin{grammarlens}[title={What you've built}]
One more word for everyday life at A2.
\end{grammarlens}

\subsection*{Your turn}
\begin{itemize}
\raggedright
  \item \emph{Say it:} \emph{jinja}
  \item \emph{Say it:} \emph{jinja}, once more
  \item \emph{Say it:} \emph{maitoshi jinja ni ikimasu} — I go to the shrine every year.
  \item \emph{Recall:} say the Japanese for to ask, then the Japanese for guidance, then say \emph{jinja} again
\end{itemize}

\begin{tcolorbox}[breakable,colback=teal!4,colframe=teal!35!black,title={Before you move on}]
What does \ja{じんじゃ} mean? (\textquotedblleft{}Shrine; Shinto shrine\textquotedblright{}.) Say it once more.
\end{tcolorbox}
\section[\texorpdfstring{\ja{てら} (tera)}{てら (tera)}]{\ja{てら} (tera) — temple; Buddhist temple}

\label{lesson:JA-C193-tera}



\begin{quote}
Before the new one: say the Japanese for guidance, then the Japanese for shrine.
\end{quote}

\subsection*{You'll want to know: \ja{てら}}
\textbf{\ja{てら}} — \emph{tera} — \textquotedblleft{}temple; Buddhist temple\textquotedblright{}.

\begin{quote}
\textbf{\ja{じんじゃの} \ja{となりに} \ja{てらが} \ja{あります。}} — \emph{jinja no tonari ni tera ga arimasu} — \textquotedblleft{}There is a temple next to the shrine.\textquotedblright{}
\end{quote}

\begin{cousinweb}
\textbf{\ja{てら}} (\emph{tera}) is the native reading (kun'yomi) given to the character for \textquotedblleft{}Buddhist temple\textquotedblright{}; its older origin is not settled (JMdict, KANJIDIC2).
\end{cousinweb}

\begin{grammarlens}[title={What you've built}]
One more word for everyday life at A2.
\end{grammarlens}

\subsection*{Your turn}
\begin{itemize}
\raggedright
  \item \emph{Say it:} \emph{tera}
  \item \emph{Say it:} \emph{tera}, once more
  \item \emph{Say it:} \emph{jinja no tonari ni tera ga arimasu} — There is a temple next to the shrine.
  \item \emph{Recall:} say the Japanese for guidance, then the Japanese for shrine, then say \emph{tera} again
\end{itemize}

\begin{tcolorbox}[breakable,colback=teal!4,colframe=teal!35!black,title={Before you move on}]
What does \ja{てら} mean? (\textquotedblleft{}Temple; Buddhist temple\textquotedblright{}.) Say it once more.
\end{tcolorbox}
\section[\texorpdfstring{\ja{えいがかん} (eigakan)}{えいがかん (eigakan)}]{\ja{えいがかん} (eigakan) — cinema; movie theatre}

\label{lesson:JA-C193-eigakan}



\begin{quote}
Before the new one: say the Japanese for shrine, then the Japanese for temple.
\end{quote}

\subsection*{You'll want to know: \ja{えいがかん}}
\textbf{\ja{えいがかん}} — \emph{eigakan} — \textquotedblleft{}cinema; movie theatre\textquotedblright{}.

\begin{quote}
\textbf{\ja{ともだちと} \ja{えいがかんに} \ja{いきます。}} — \emph{tomodachi to eigakan ni ikimasu} — \textquotedblleft{}I go to the cinema with a friend.\textquotedblright{}
\end{quote}

\begin{cousinweb}
\textbf{\ja{えいがかん}} (\emph{eigakan}) is Sino-Japanese (kango): three Chinese characters read with their on'yomi, readings borrowed from Chinese. They are \emph{ei}, \textquotedblleft{}reflect\textquotedblright{} (Middle Chinese \emph{'jaengH}), \emph{ga}, \textquotedblleft{}picture\textquotedblright{} (\emph{hweaH}) and \emph{kan}, \textquotedblleft{}building\textquotedblright{} (\emph{kwanH}) (KANJIDIC2; Middle Chinese in Baxter's transcription).
\end{cousinweb}

\begin{grammarlens}[title={What you've built}]
One more word for everyday life at A2.
\end{grammarlens}

\subsection*{Your turn}
\begin{itemize}
\raggedright
  \item \emph{Say it:} \emph{eigakan}
  \item \emph{Say it:} \emph{eigakan}, once more
  \item \emph{Say it:} \emph{tomodachi to eigakan ni ikimasu} — I go to the cinema with a friend.
  \item \emph{Recall:} say the Japanese for shrine, then the Japanese for temple, then say \emph{eigakan} again
\end{itemize}

\begin{tcolorbox}[breakable,colback=teal!4,colframe=teal!35!black,title={Before you move on}]
What does \ja{えいがかん} mean? (\textquotedblleft{}Cinema; movie theatre\textquotedblright{}.) Say it once more.
\end{tcolorbox}
\section[\texorpdfstring{\ja{びじゅつかん} (bijutsukan)}{びじゅつかん (bijutsukan)}]{\ja{びじゅつかん} (bijutsukan) — art museum; art gallery}

\label{lesson:JA-C193-bijutsukan}



\begin{quote}
Before the new one: say the Japanese for temple, then the Japanese for cinema.
\end{quote}

\subsection*{You'll want to know: \ja{びじゅつかん}}
\textbf{\ja{びじゅつかん}} — \emph{bijutsukan} — \textquotedblleft{}art museum; art gallery\textquotedblright{}.

\begin{quote}
\textbf{\ja{しゅうまつ} \ja{びじゅつかんに} \ja{いきます。}} — \emph{shūmatsu bijutsukan ni ikimasu} — \textquotedblleft{}I go to the art museum at the weekend.\textquotedblright{}
\end{quote}

\begin{cousinweb}
\textbf{\ja{びじゅつかん}} (\emph{bijutsukan}) is Sino-Japanese (kango): three Chinese characters read with their on'yomi, readings borrowed from Chinese. They are \emph{bi}, \textquotedblleft{}beauty\textquotedblright{} (Middle Chinese \emph{mijX}), \emph{jutsu}, \textquotedblleft{}art, skill\textquotedblright{} (\emph{zywit}) and \emph{kan}, \textquotedblleft{}building\textquotedblright{} (\emph{kwanH}) (KANJIDIC2; Middle Chinese in Baxter's transcription).
\end{cousinweb}

\begin{grammarlens}[title={What you've built}]
One more word for everyday life at A2.
\end{grammarlens}

\subsection*{Your turn}
\begin{itemize}
\raggedright
  \item \emph{Say it:} \emph{bijutsukan}
  \item \emph{Say it:} \emph{bijutsukan}, once more
  \item \emph{Say it:} \emph{shūmatsu bijutsukan ni ikimasu} — I go to the art museum at the weekend.
  \item \emph{Recall:} say the Japanese for temple, then the Japanese for cinema, then say \emph{bijutsukan} again
\end{itemize}

\begin{tcolorbox}[breakable,colback=teal!4,colframe=teal!35!black,title={Before you move on}]
What does \ja{びじゅつかん} mean? (\textquotedblleft{}Art museum; art gallery\textquotedblright{}.) Say it once more.
\end{tcolorbox}
\section[\texorpdfstring{\ja{みえる} (mieru)}{みえる (mieru)}]{\ja{みえる} (mieru) — to be visible; to be seen, to seem}

\label{lesson:JA-C193-mieru}



\begin{quote}
Before the new one: say the Japanese for cinema, then the Japanese for art museum.
\end{quote}

\subsection*{You'll want to know: \ja{みえる}}
\textbf{\ja{みえる}} — \emph{mieru} — \textquotedblleft{}to be visible; to be seen, to seem\textquotedblright{}.

\begin{quote}
\textbf{\ja{あの} \ja{たかい} \ja{たてものが} \ja{みえます。}} — \emph{ano takai tatemono ga miemasu} — \textquotedblleft{}That tall building can be seen.\textquotedblright{}
\end{quote}

Polite \textbf{\ja{みえます}}, \emph{miemasu}; an ichidan verb.

\begin{cousinweb}
\textbf{\ja{みえる}} (\emph{mieru}) is a native Japanese word, inherited rather than borrowed, built on \emph{mi}, \textquotedblleft{}seeing\textquotedblright{}; its kanji was attached later as a kun'yomi (JMdict, KANJIDIC2).
\end{cousinweb}

\begin{grammarlens}[title={What you've built}]
One more word for everyday life at A2.
\end{grammarlens}

\subsection*{Your turn}
\begin{itemize}
\raggedright
  \item \emph{Say it:} \emph{mieru}
  \item \emph{Say it:} \emph{mieru}, once more
  \item \emph{Say it:} \emph{ano takai tatemono ga miemasu} — That tall building can be seen.
  \item \emph{Recall:} say the Japanese for cinema, then the Japanese for art museum, then say \emph{mieru} again
\end{itemize}

\begin{tcolorbox}[breakable,colback=teal!4,colframe=teal!35!black,title={Before you move on}]
What does \ja{みえる} mean? (\textquotedblleft{}To be visible; to be seen, to seem\textquotedblright{}.) Say it once more.
\end{tcolorbox}
